% Part: first-order-logic
% Chapter: syntax-and-semantics
% Section: formation-sequences

\documentclass[../../../include/open-logic-section]{subfiles}

\begin{document}

\olfileid{pl}{syn}{fseq}

\olsection{रचनाक्रमिका}

!!{formula}s ची विगमनाधारित व्याख्या देणे आणि त्याला पूरक तंत्र म्हणून
!!{formula}s चे गुणधर्म विगमनाने सिद्ध करणे हा सुबक व कार्यक्षम
दृष्टिकोन आहे. तथापि, !!{formula}s च्या रचनेचा तळापासून वर जाणारा,
टप्प्याटप्प्याचा दृष्टिकोनही उपयुक्त ठरू शकतो; येथे आपण तो
\emph{रचनाक्रमिका} या संकल्पनेच्या साहाय्याने मांडतो.

\begin{defn}[सूत्रांसाठी रचनाक्रमिका]
\ollabel{defn:fseq-frm}
भाषा~$\Lang L_0$ मधील चिन्हांपासून बनलेल्या चिन्हमालांची सांत क्रमिका
$\tuple{!A_0,\dotsc,!A_n}$ ही $!A$ साठी \emph{रचनाक्रमिका} असते, जर
$!A \ident !A_n$ आणि प्रत्येक $i \leq n$ साठी एकतर $!A_i$ हे आण्विक
सूत्र असेल किंवा पुढीलपैकी एखादी अट पूर्ण करणारे $j,k < i$ अस्तित्वात
असतील:
\begin{enumerate}
    \tagitem{prvNot}{$!A_i \ident \lnot !A_j$.}{}%
    \tagitem{prvAnd}{$!A_i \ident (!A_j \land !A_k)$.}{}%
    \tagitem{prvOr}{$!A_i \ident (!A_j \lor !A_k)$.}{}%
    \tagitem{prvIf}{$!A_i \ident (!A_j \lif !A_k)$.}{}%
    \tagitem{prvIff}{$!A_i \ident (!A_j \liff !A_k)$.}{}%
\end{enumerate}
\end{defn}

\begin{ex}
\[
\tuple{
    \Obj p_0,
    \Obj p_1,
    (\Obj p_1 \land \Obj p_0),
    \lnot (\Obj p_1 \land \Obj p_0)
}
\]
ही $\lnot (\Obj p_1 \land \Obj p_0)$ ची एक रचनाक्रमिका आहे; तसेच
पुढील क्रमिकाही तिची रचनाक्रमिका आहे:
\[
\tuple{
    \Obj p_0,
    \Obj p_1,
    \Obj p_0,
    (\Obj p_1 \land \Obj p_0),
    (\Obj p_0 \lif \Obj p_1),
    \lnot (\Obj p_1 \land \Obj p_0)
}.
\]
%
दुसऱ्या उदाहरणावरून दिसते की रचनाक्रमिकेत `अडगळ' असू शकते: अनावश्यक
असलेली किंवा रचनेत काहीही हातभार न लावणारी सूत्रे.
\end{ex}

\begin{prop}
\ollabel{prop:formed}
$\Frm[L_0]$ मधील प्रत्येक !!{formula}~$!A$ ला एक रचनाक्रमिका असते.
\end{prop}

\begin{proof}
$!A$ आण्विक आहे असे समजा. मग क्रमिका~$\tuple{!A}$ ही $!A$ साठी
रचनाक्रमिका आहे.
%
आता $!B$ आणि~$!C$ यांना अनुक्रमे $\tuple{!B_0,\dotsc,!B_n}$ आणि
$\tuple{!C_0,\dotsc,!C_m}$ या रचनाक्रमिका आहेत असे समजा.
%
\begin{enumerate}
  \tagitem{prvNot}{जर $!A \ident \lnot !B$, तर
      $\tuple{!B_0,\dotsc,!B_n,\lnot !B_n}$ ही $!A$ साठी
      रचनाक्रमिका आहे.}{}
  \tagitem{prvAnd}{जर $!A \ident (!B \land !C)$, तर
      $\tuple{!B_0,\dotsc,!B_n,!C_0,\dotsc,!C_m,(!B_n \land !C_m)}$
      ही $!A$ साठी रचनाक्रमिका आहे.}{}
  \tagitem{prvOr}{जर $!A \ident (!B \lor !C)$, तर
      $\tuple{!B_0,\dotsc,!B_n,!C_0,\dotsc,!C_m,(!B_n \lor !C_m)}$
      ही $!A$ साठी रचनाक्रमिका आहे.}{}
  \tagitem{prvIf}{जर $!A \ident (!B \lif !C)$, तर
      $\tuple{!B_0,\dotsc,!B_n,!C_0,\dotsc,!C_m,(!B_n \lif !C_m)}$
      ही $!A$ साठी रचनाक्रमिका आहे.}{}
  \tagitem{prvIff}{जर $!A \ident (!B \liff !C)$, तर
      $\tuple{!B_0,\dotsc,!B_n,!C_0,\dotsc,!C_m,(!B_n \liff !C_m)}$
      ही $!A$ साठी रचनाक्रमिका आहे.}{}
\end{enumerate}
!!{formula}s वरील विगमनाच्या तत्त्वानुसार प्रत्येक !!{formula} ला एक
रचनाक्रमिका असते.
\end{proof}

याचे उलट विधानही आपण सिद्ध करू शकतो. ते महत्त्वाचे आहे, कारण त्यामुळे
सूत्रे परिभाषित करण्याच्या आपल्या दोन्ही पद्धती सममूल्य आहेत, म्हणजे
त्या समान परिणाम देतात, हे दिसते. तसेच रचनाक्रमिकांच्या लांबीवर साधे
विगमन वापरून सूत्रांविषयीची प्रमेये सिद्ध करता येतात.

\begin{lem}
\ollabel{lem:fseq-init}
$\tuple{!A_0,\dotsc,!A_n}$ ही $!A_n$ साठी रचनाक्रमिका आहे आणि
$k \leq n$ असे समजा. मग $\tuple{!A_0,\dotsc,!A_k}$ ही $!A_k$ साठी
रचनाक्रमिका आहे.
\end{lem}

\begin{proof}
अभ्यासप्रश्न.
\end{proof}

\begin{thm}
\ollabel{thm:fseq-frm-equiv}
$\Frm[L_0]$ हा भाषा~$\Lang L_0$ मधील ज्या सर्व चिन्हमालांना
रचनाक्रमिका आहे, त्यांचा संच आहे.
\end{thm}

\begin{proof}
भाषा~$\Lang L_0$ मधील ज्या सर्व चिन्हमालांना रचनाक्रमिका आहे, त्यांचा
संच~$F$ असू द्या. $\Frm[L_0] \subseteq F$ हे
\olref[pl][syn][fseq]{prop:formed} मध्ये आपण पाहिले आहे; म्हणून आता उलट
समावेश सिद्ध करू.

$!A$ ला $\tuple{!A_0,\dotsc,!A_n}$ ही रचनाक्रमिका आहे असे समजा.
$n$ वरील प्रबल विगमनाने $!A \in \Frm[L_0]$ हे सिद्ध करू. लांबी
$m < n$ असलेली रचनाक्रमिका ज्या प्रत्येक चिन्हमालेला आहे, ती
$\Frm[L_0]$ मध्ये आहे, हे आपले विगमन गृहीतक आहे. रचनाक्रमिकेच्या
व्याख्येनुसार एकतर $!A_n$ आण्विक आहे किंवा पुढीलपैकी एखादी अट पूर्ण
करणारे $j,k < n$ अस्तित्वात असले पाहिजेत:
\begin{enumerate}
    \tagitem{prvNot}{$!A_n \ident \lnot !A_j$.}{}%
    \tagitem{prvAnd}{$!A_n \ident (!A_j \land !A_k)$.}{}%
    \tagitem{prvOr}{$!A_n \ident (!A_j \lor !A_k)$.}{}%
    \tagitem{prvIf}{$!A_n \ident (!A_j \lif !A_k)$.}{}%
    \tagitem{prvIff}{$!A_n \ident (!A_j \liff !A_k)$.}{}%
\end{enumerate}
आता प्रत्येक शक्यतेचा स्वतंत्र विचार करू. $!A_n$ आण्विक असेल, तर
$!A_n \in \Frm[L_0]$. त्याऐवजी $!A \ident
(!A_j \land !A_k)$ असे समजा.
%READERNOTE{MRPL-002: गोठवलेल्या इंग्रजीत या एकाच शक्यतेत विन्यासात्मक अनन्यता-चिन्हाऐवजी पर्यायी सममूल्यता-चिन्ह आले आहे. आधीची स्पष्ट व्याख्या आणि इतर सर्व रचना-बाबी यांच्याशी जुळण्यासाठी येथे विन्यासात्मक अनन्यता वापरली आहे.}
\olref[pl][syn][fseq]{lem:fseq-init} नुसार
$\tuple{!A_0,\dotsc,!A_j}$ आणि $\tuple{!A_0,\dotsc,!A_k}$ या अनुक्रमे
$!A_j$ आणि~$!A_k$ साठी रचनाक्रमिका आहेत. या $!A$ च्या
रचनाक्रमिकेच्या उचित आरंभीच्या उपक्रमिका असल्यामुळे दोन्हींची लांबी
$n$ पेक्षा कमी आहे. म्हणून विगमन गृहीतकानुसार $!A_j$ आणि~$!A_k$ ही
$\Frm[L_0]$ मध्ये आहेत; त्यामुळे !!a{formula} च्या व्याख्येनुसार
$(!A_j \land !A_k)$ हेही त्यात आहे. इतर शक्यता समांतर युक्तिवादाने
सिद्ध होतात.
\end{proof}

\end{document}
